\documentclass[11pt]{article}
\usepackage[margin=1in]{geometry}
\usepackage{amsmath,amssymb,amsthm}
\usepackage{graphicx}
\usepackage[hidelinks]{hyperref}

\newtheorem{theorem}{Theorem}
\newtheorem{lemma}{Lemma}
\newtheorem{remark}{Remark}

\newcommand{\UMD}{\mathrm{UMD}}
\newcommand{\E}{\mathbb{E}}
\newcommand{\ellone}{\ell_1}

\title{A No-Go Result for the Standard \(L^1\) Decoupling Examples in the Type 2 Non-UMD Question}
\author{run fa\_banach\_001, agent\_lane\_15}
\date{2026-06-28}

\begin{document}
\maketitle

\section*{Status}

This is a candidate partial result.  It does not solve the full question in
Cox--Gorajski.  It proves that two standard \(L^1\)-based sources of non-UMD
spaces with the decoupling property cannot produce a type 2 non-UMD example.

\section*{Source Metadata}

The source paper is Sonja Cox and Mariusz Gorajski,
\emph{Vector-valued stochastic delay equations -- a semigroup approach},
arXiv:1006.5349.  The local source TeX line is
\path{data/parsed/arxiv_sources/1006.5349/source.tex:541}.  The rendered source
evidence below is page 12 of the locally compiled source PDF.

\begin{center}
\includegraphics[width=\linewidth]{figures/open_problem_crop.png}
\end{center}

\section*{Source Question}

In Section 4.1, Cox and Gorajski assume that their stochastic delay equation is
set in a type 2 UMD Banach space \(E\).  They add the following parenthetical
comment:
\begin{quote}
The results in this article are also valid if \(E\) is a type 2 Banach space
with the decoupling property but we do not know of any such spaces that are not
in fact UMD spaces.
\end{quote}

\section*{Definitions and Background}

A Banach space has Gaussian type 2 if there is a constant \(T\) such that for
all finite choices \(x_1,\ldots,x_n\),
\[
   \left(\E\left\|\sum_{j=1}^n \gamma_j x_j\right\|^2\right)^{1/2}
   \leq T\left(\sum_{j=1}^n \|x_j\|^2\right)^{1/2},
\]
where the \(\gamma_j\)'s are independent standard Gaussian variables.

We use the following standard facts.
\begin{enumerate}
  \item Scalar \(L^1(\mu)\) has cotype 2, and cotype 2 passes to subspaces and
  isomorphs.
  \item Kwapien's theorem says that a Banach space with type 2 and cotype 2 is
  isomorphic to a Hilbert space.
  \item Hilbert spaces are UMD, and finite direct sums of UMD spaces are UMD.
  \item Cox--Veraar proved that \(L^1(S;Y)\) has the decoupling property when
  \(S\) is sigma-finite and \(Y\) is UMD.
\end{enumerate}

\section*{Proof Intuition}

The known non-UMD decoupling examples are \(L^1\)-flavored.  But type 2 is
hostile to \(\ell_1^n\)-geometry.  In scalar \(L^1\) subspaces the obstruction
is even stronger: scalar \(L^1\) has cotype 2, so adding type 2 collapses the
space to Hilbert by Kwapien's theorem.  In the Bochner examples \(L^1(S;Y)\),
an infinite or nonatomic measure algebra gives isometric copies of
\(\ell_1^n\) for arbitrarily large \(n\), contradicting type 2.  Thus the only
type 2 Bochner cases are finite \(L^1\)-sums of the original UMD space \(Y\).

\section*{Main Partial Result}

\begin{theorem}[No-go for the standard \(L^1\) decoupling examples]
Let \(E\) be a Banach space belonging to either of the following two classes.
\begin{enumerate}
  \item \(E\) is isomorphic to a closed subspace of a scalar \(L^1(\mu)\).
  \item \(E=L^1(S;Y)\), where \((S,\Sigma,\mu)\) is sigma-finite and \(Y\) is
  UMD.
\end{enumerate}
If \(E\) has type 2, then \(E\) is UMD.  More precisely, in case (1), \(E\) is
isomorphic to a Hilbert space.  In case (2), either \(Y=\{0\}\), or the measure
algebra of \(S\) has only finitely many atoms and \(E\) is isomorphic to a
finite \(\ell_1\)-sum of copies of \(Y\).
\end{theorem}

\begin{proof}
Assume first that \(E\) is isomorphic to a closed subspace of a scalar
\(L^1(\mu)\).  Since \(L^1(\mu)\) has cotype 2 and cotype 2 passes to closed
subspaces and to isomorphic copies, \(E\) has cotype 2.  By hypothesis \(E\)
also has type 2.  Kwapien's theorem therefore implies that \(E\) is isomorphic
to a Hilbert space.  Hence \(E\) is UMD.

Now assume that \(E=L^1(S;Y)\), where \(S\) is sigma-finite and \(Y\) is UMD.
If \(Y=\{0\}\), there is nothing to prove, so choose \(y_0\in Y\) with
\(\|y_0\|=1\).  Suppose that \(S\) contains \(n\) pairwise disjoint measurable
sets \(A_1,\ldots,A_n\), each of finite positive measure.  Then
\[
   (a_1,\ldots,a_n)\mapsto
   \sum_{k=1}^n a_k\,\mu(A_k)^{-1}\mathbf{1}_{A_k}y_0
\]
is an isometric embedding of \(\ell_1^n\) into \(L^1(S;Y)\).  Indeed, the
supports are disjoint, so the \(L^1(S;Y)\)-norm of the image is
\(\sum_{k=1}^n |a_k|\).

If such disjoint families existed for arbitrarily large \(n\), then
\(L^1(S;Y)\) would contain isometric copies of \(\ell_1^n\) for arbitrarily
large \(n\).  This is incompatible with type 2: applying the type 2 inequality
inside \(\ell_1^n\) to the coordinate vectors \(e_1,\ldots,e_n\) gives
\[
  \left(\E\left(\sum_{k=1}^n |\gamma_k|\right)^2\right)^{1/2}
  \leq T\sqrt{n},
\]
whereas the left side is at least
\(n\,\E|\gamma_1|=\sqrt{2/\pi}\,n\).  Thus the type 2 constant would have to
grow at least like \(\sqrt n\), a contradiction.

Therefore a type 2 space \(L^1(S;Y)\) admits a uniform bound on the size of a
family of pairwise disjoint positive finite-measure sets.  Since \(S\) is
sigma-finite, this forces the measure algebra to be finite atomic: a nonatomic
positive-measure part can be split into arbitrarily many positive finite
measure pieces, and infinitely many atoms also give arbitrarily large disjoint
families.  Hence \(L^1(S;Y)\) is isomorphic, indeed after rescaling atom
weights isometric, to a finite \(\ell_1\)-direct sum of copies of \(Y\).
Finite sums preserve the UMD property, so \(E\) is UMD.
\end{proof}

\section*{Consequence for the Cox--Gorajski Aside}

The theorem rules out the most visible \(L^1\)-based non-UMD decoupling routes
as answers to the Cox--Gorajski parenthetical question.  In particular, the
non-UMD examples \(L^1(S;Y)\) from Cox--Veraar cannot simultaneously have type
2 unless the measure part is finite atomic, in which case they are UMD.
Likewise, scalar \(L^1\)-subspace examples with type 2 are Hilbertian.

\section*{Limitations and Novelty Check}

This is only a partial no-go result.  It does not prove that every type 2 space
with the decoupling property is UMD.  It also does not address hypothetical
decoupling examples outside the scalar \(L^1\)-subspace and Cox--Veraar
Bochner \(L^1(S;Y)\) families.

The lightweight run indexes were searched for arXiv:1006.5349, the source
title, and the terms \texttt{type 2}, \texttt{decoupling property},
\texttt{UMD}, \texttt{L1}, \texttt{Cox Veraar}, and related local source-corpus
keywords.  No duplicate packet for this exact partial no-go theorem was found.

\section*{References}

\begin{thebibliography}{9}
\bibitem{CG}
S. Cox and M. Gorajski,
\emph{Vector-valued stochastic delay equations -- a semigroup approach},
arXiv:1006.5349.

\bibitem{CV}
S. Cox and M. Veraar,
\emph{Some remarks on tangent martingale difference sequences in
\(L^1\)-spaces}, arXiv:0801.0695.

\bibitem{KLW}
N. J. Kalton, E. Lorist, and L. Weis,
\emph{Euclidean structures and operator theory in Banach spaces},
arXiv:1912.09347.  See the discussion of the Kwapien--Maurey theorem and the
type 2 plus cotype 2 Hilbertian consequence.

\bibitem{AK}
F. Albiac and N. J. Kalton,
\emph{Topics in Banach Space Theory}, second edition, Springer, 2016.
\end{thebibliography}

\end{document}
